% !TEX program = xelatex
% 编译：xelatex BiT128_knowledge_primer.tex （编译两遍以生成交叉引用）
\documentclass[11pt,a4paper]{ctexart}

\usepackage[margin=2.3cm]{geometry}
\usepackage{amsmath,amssymb}
\usepackage{booktabs}
\usepackage{longtable}
\usepackage[colorlinks=true,linkcolor=blue!55!black,urlcolor=blue!55!black]{hyperref}
\usepackage{tcolorbox}
\tcbuselibrary{skins,breakable}
\newcommand{\Z}{\mathbb{Z}}
\newcommand{\Zq}{\mathbb{Z}_q}
\usepackage{fancyhdr}

\pagestyle{fancy}
\fancyhf{}
\fancyhead[L]{\small BiT-128 攻击知识脉络（入门版）}
\fancyhead[R]{\small\thepage}

\newtcolorbox{intuition}[1][]
  {colback=blue!4,colframe=blue!55!black,title={直觉 #1},fonttitle=\bfseries,breakable}
\newtcolorbox{checkpoint}[1][]
  {colback=green!4,colframe=green!45!black,title={检查点 #1},fonttitle=\bfseries,breakable}
\newtcolorbox{warning}[1][]
  {colback=red!4,colframe=red!55!black,title={容易想错的地方 #1},fonttitle=\bfseries,breakable}
\newtcolorbox{story}[1][]
  {colback=orange!5,colframe=orange!60!black,title={实验室日志 #1},fonttitle=\bfseries,breakable}

\title{\bfseries 从零开始理解 BiT-128 攻击\\[0.3em]
\large ——写给只了解皮毛的同学的完整知识脉络（v2，含 9-28 更正）}
\author{GLM-5.3 Flash max \and Jy L\\{\small 配套《BiT-128 签名方案的公开数据统计攻击》技术报告使用}}
\date{2026 年 9 月}

\begin{document}
\maketitle

\begin{tcolorbox}[colback=blue!4,colframe=blue!55!black,title=这份文档怎么读,fonttitle=\bfseries]
\textbf{读者假设}：你会模运算（$17 \equiv 4 \pmod{13}$ 这种）、知道哈希函数是"单向压碎机"、
知道公钥私钥的大概概念。除此之外\textbf{不需要}任何格密码、拒绝采样或 NGCC 的背景——
所有需要的东西都在这里从头建立。

\textbf{结构}：第 1--4 章打地基（数学工具、签名骨架、拒绝采样、BiT 流程）；
第 5--9 章是攻击本身（每章回答一个问题）；第 10 章是实验时间线；
第 11 章是交叉验证；第 12 章是 FAQ。第 13 章有符号速查表。

\textbf{配套}：严谨版是同目录的 \texttt{BiT128\_attack\_report.pdf}（技术报告）；
演讲版是 \texttt{BiT128\_seminar\_beamers.pdf}（讨论班）。
\end{tcolorbox}

\tableofcontents
\newpage

\section{数学地基：环与"会环绕的多项式"}

\subsection{模运算：一根 26881 刻度的圆尺}

$\Z_q$ 就是"把所有整数按 $q=26881$ 收卷成一个圆"：$17 \equiv 17$，
$26881 \equiv 0$，$53695 \equiv -67$。习惯上每个数用它在 $[-q/2, q/2)$ 里的
\textbf{中心代表} $\mathrm{center}(x)$ 表示。$-1$ 和 $q-1$ 是同一个东西的两件外衣。

\begin{intuition}
把 $\Z_q$ 想成一根 26881 刻度的圆尺弯成一个环。"加一个大数"等于在环上转了很多圈
又回到附近；"模 $q$"就是只记你最终停在哪一刻度。
\end{intuition}

\subsection{多项式：256 个数字排成一列，乘法会"环绕"}

格密码对\textbf{系数很小的多项式}做运算：
$f(x) = f_0 + f_1 x + \cdots + f_{255} x^{255}$，$f_i \in \Z_q$。
乘法规定 $x^{256} \equiv -1$：最高次项乘上去会\textbf{翻个符号绕回常数项}
——负循环（negacyclic）。

\begin{intuition}[小例子：$d=4$]
$(a_0 + a_1 x + a_2 x^2 + a_3 x^3)\cdot(b_0 + b_1 x)$ 的普通乘积里有一项
$a_3 b_1 x^4$；因为 $x^4 \equiv -1$，它变成 $-a_3b_1$ 加回常数项。一般地：
\[
  (f*g)[j] = \sum_{k} f[k]\, g[(j-k) \bmod d]\cdot (-1)^{\mathbf{1}[j<k]} ,
\]
"下标绕回去时配一个负号"。这个负号因子是双刃剑：它让环代数成立，
也让攻击里的对齐统计要小心符号——第 7 章反复用到。
\end{intuition}

\begin{checkpoint}
$d=4$ 时计算 $(1+x)*(1-x)$。答案：$1-x^2$。再乘 $x^2$ 呢？$-1+x^2$
——$x^2$ 项翻了符号（$x^4\equiv-1$）。
\end{checkpoint}

\subsection{"小乘小、加起来不大"}

BiT 所有秘密量都很小：$s_0, e$ 系数在 $\{-1,0,1\}$；掩码 $y$ 系数在 $\pm2047$ 内。
30 项 $\pm1$ 相加落在 $[-30,30]$。\textbf{"小秘密卷上已知的小东西，得到的还是
（相对）小的、分布可计算的量"}——全文反复出现的主题。

\section{数字签名的骨架与 Fiat--Shamir}

\subsection{签名要证明什么}

数字签名：持有私钥的人能对任意消息 $m$ 产生凭证 $\sigma$，任何人可验证，
但看过再多签名也\textbf{造不出}新消息的签名（EUF-CMA）。格签名是零知识证明的变体：
证明"我知道一个很短的 $s_0$ 满足 $A_0 s_0 \equiv$ 某公开值"，但不泄露 $s_0$。

\subsection{Fiat--Shamir：把"问答"压成"独白"}

原始零知识证明是对话：证明者扔掩码 $y$、给承诺 $w$；验证者随机出题 $c$；
证明者回答 $z = y + (\text{含 } s_0 \text{ 的项})$。因为 $y$ 均匀，$z$ 看起来均匀。
\textbf{Fiat--Shamir} 把考题换成 $\tilde c = H(\text{消息}\|w)$——哈希充当
"懒惰但公正的考官"。BiT 正是这样。

\subsection{拒绝采样：为什么"扔了重扔"}

$z = y + \varepsilon(s_0*c)$ 里混着秘密项，$z$ 的分布会向 $\pm s_0*c$ 偏一点
——统计足够多 $z$ 就能恢复秘密。补救：\textbf{算完检查，不像标准分布就整个扔掉重来}。
发布后的 $z$ 都来自标准分布。

\begin{warning}
拒绝采样堵的是"$z$ 的\textbf{边缘分布}失真"。它\textbf{不}保证"$z$ 与扔硬币结果
的\textbf{联合}分布无关"。BiT 的缺口正在后者——第 5 章。
\end{warning}

\section{双峰技巧与共享符号位}

\subsection{双峰：掩码减半的代价}

掩码要够宽才能盖住秘密项。双峰观察：挑战可整体乘 $\pm1$ 时，响应的两个候选分布
对称——掩码可\textbf{减半}。签名时掷一个比特 $b$，$\varepsilon=(-1)^b$，
$c_{\mathrm{signed}} = \varepsilon c_{\mathrm{can}}$——\textbf{整条签名共享同一个 $\varepsilon$}。

\subsection{这个比特就是猎物}

两个事实：
(1) $\varepsilon$ 是\textbf{一个}比特，但影响整条签名的\textbf{所有}分量
（$z_0,z_1,z_2,h$ 全含其符号）；
(2) 规范 §3.1.3 声称"$z$ 与 $b$ 独立"——按字面理解攻击者无法知道 $b$。
第 5 章会看到：每个分量\textbf{单独}看或许独立，合起来不是。

\section{BiT-128 签名全流程（带直觉重走）}

参数：$d=256$，$q=26881$，$\tau=30$，$\gamma_2=682$，$K=L=3$。
秘密 $s_0\in\{-1,0,1\}^{3\times256}$，$e$ 同型；$b = A_0 s_0 + e = 16b_1 + b_0$
（$b_1$ 公开、$b_0\in[-8,8)$ 保密）。

\subsection{三份掩码，三种用途}

\begin{center}
\begin{tabular}{llll}
\toprule
掩码 & 宽度 & 参与的响应 & 秘密搭档 \\
\midrule
$y_0$ & $\pm7$（3-bit 三角） & $z_0 = y_0 + c_{\mathrm{signed}}$ & 无（$c$ 是公开形状）\\
$y_1$ & $\pm1023$（10-bit） & $z_1 = y_1 + s_0*c_{\mathrm{signed}}$ & $s_0$ \\
$y_2$ & $\pm2047$（11-bit） & $z_2 = y_2 + e*c_{\mathrm{signed}}$ & $e$ \\
\bottomrule
\end{tabular}
\end{center}

$z_0$ 特殊：秘密部分只是 $\pm c_{\mathrm{can}}$ 形状的公开挑战，掩码可极窄。

\subsection{挑战从哪来：承诺再提问}

每轮：$w = A_0 y_1 - 16 b_1 y_0 + \hat q y_0$（行 0）$+ y_2$，其高位摘要进哈希：
$\tilde c = H(\mu\|w_1)$，展开出挑战 $c_{\mathrm{can}}$（30 个 $+1$）。
\textbf{掩码先承诺，挑战后到}。

\subsection{三个拒绝阶段}

(1) $z_0,z_1$ 检查（失败只重扔 $y_0,y_1$）；(2) $z_2$ 检查（失败全扔）；
(3) hint 与范数检查（失败全扔）。注意：\textbf{"接受"本身是以秘密为条件的事件}。

\subsection{hint：给验证器的小抄}

验证器重算 $w'$，想知道原来的高位 $w_1$；但 $w'$ 与签名时的 $w$ 差一点，高位可能差一两格。
$h$ 就是\textbf{纠错小抄}：告诉验证器每坐标的格子号怎么调。调完哈希对上，验证通过。

\subsection{一条代数金线}

密钥生成保证
\begin{equation}
A_0 s_0 \equiv 16 b_1 + (b_0 - e) \pmod q,
\qquad R := b_0 - e \in [-9,8]^{3\times256}.
\label{eq:primer-R}
\end{equation}
$R$ 公钥可计算；反过来知道 $R$ 也能解出 $s_0$（第 9 章）。\textbf{$R$ 是连接公开世界
与秘密世界的桥}——攻击目标从"恢复 $s_0$"等价变成"恢复 $R$"。

\section{缺口在哪里：边缘独立 $\ne$ 联合独立}

\subsection{玩具模型}

两个坐标的玩具签名：拒绝规则说"$|z|$ 太大就扔"。\textbf{被发布}的签名里两坐标的
联合分布与理想模型不同——每坐标边缘可能合规，但\textbf{相关性}泄露秘密。

\begin{intuition}
海关只检查每件行李是否超标；但"行李超标的旅客往往有同伙"——联合检查（看相关性）
能发现走私网络。拒绝采样是海关，联合分布是走私网络。
\end{intuition}

\subsection{BiT 的具体情况}

规范 §3.1.3 声称独立。而联合响应实际保留秘密相关的跨坐标相关——我们自己的实验从数据上证实并精确刻画（第 7 章）：
同键双数据集的统计量零位置模式相关达 $0.9907$——确定性的、随 $N$ 增长的指纹。

\section{通道一：$z_0$——先偷到那个比特}

支撑位置上 $z_0 = y_0 \pm 1$：$\varepsilon=+1$ 发布值集中在 $[-6,7]$（右截断），
$\varepsilon=-1$ 集中在 $[-7,6]$（左截断）。不对称 $\Rightarrow$ 逐位置对数似然比
$\ln[f_0(z-1)/f_0(z+1)]$ 求和得 $\Lambda$，$\hat\varepsilon=\mathrm{sign}(\Lambda)$。
$10^6$ 真实签名实测 $96.19\%$（边界事件覆盖 62.3\%、逐事件 100\%）。
\textbf{攻击者现在知道每条签名的 $\varepsilon$（96.2\% 置信）}。

关于 $z_1$ 通道：它携带 $s_0$ 的一阶信息。我们早期用压缩统计量
$\psi(z)=\mathrm{sign}(z)/(1024-|z|)$ 加简单阈值估计它时观察到不收敛，于是转向 hint 通道。
【2026-09-28 更正】不收敛是那个压缩统计量与阈值判决的缺陷——$z_1$ 通道用正确的
加权回归是收敛的（我们后来完整复现了这条路线，两种路线恢复出\textbf{同一把}等价私钥）。
hint 通道的价值在于：它给出 $R$ 的\textbf{直接}一阶观测、确定性结构可写成精确公式并求逆。

\section{通道二：hint——$y_2$ 的精确抵消}

\subsection{第一步：写出 $h$ 的量测对象}

$h[j] = \mathrm{fold}_{40}(\mathrm{cell}(w[j]) - \mathrm{cell}(w_a[j]))$，
$\mathrm{cell}$ 把环切成 40 格（格宽 $\gamma_2=682$），
$w_a[j] = w[j] - z_2[j] + (b_0*c_{\mathrm{signed}})[j]$（行 0 且 $b=0$ 时再加
$c_{\mathrm{signed}}[j]$）。

\subsection{第二步：展开差值，$y_2$ 消失}

\[
w[j]-w_a[j] = z_2[j] - (b_0*c)[j] - \cdots
= y_2[j] + (e*c)[j] - (b_0*c)[j] - \cdots
= y_2[j] - (R*c_{\mathrm{signed}})[j] - \cdots
\]

\begin{intuition}[魔法的本质]
$w$ 里有 $+y_2$，$z_2$ 里也有 $+y_2$。签名者把 $z_2$ \textbf{减}回去——
$y_2$ 一进一出\textbf{精确抵消}。设计者想要"hint 只依赖好处理的量"；
没注意抵消之后\textbf{剩下的是 $R$}。
\end{intuition}

\subsection{第三步：读出模型}

$\Delta[j] = y_2[j] - a[j]$（$a[j] = (R*c_{\mathrm{signed}})[j]$ 加行 0 项），
$\mathbb{E}[h[j]|\Delta[j]] \approx \Delta[j]/\gamma_2$：$h$ 是 $\Delta/\gamma_2$ 的
7 级量化读数。MC 斜率 $-1.4826\times10^{-3}$ vs 理论 $-1.4663\times10^{-3}$。

\begin{checkpoint}
为什么 $y_2$ 是噪声而非信号？——$y_2$ 每条签名全新、与密钥无关；
$a[j]$ 里藏着\textbf{固定的} $R$。平均会留下 $a$、洗掉 $y_2$。
\end{checkpoint}

\subsection{第四步：对齐累加}

\[
\Phi_R[i][m] = \sum_{n,t} \hat\varepsilon_n\,\mathrm{flip}\cdot
h_n[(m+p_{n,t}) \bmod d]
\]
$\hat\varepsilon_n$（通道一的标签）把每条签名的信号对齐同号相加。
$N$ 条签名 $\times\ \tau=30$ 位置 $= N\tau$ 个一阶观测，全部指向同一个 $R$。

\section{场定律：剩下那个"洗不掉"的东西}

\subsection{意外：直接累加失败}

$N=10^6$ 累加后直接分类失败：残差是独立噪声预期的 23--37 倍，同键另一套数据在
同样位置呈同样模式（相关 0.99）。像被"染色"——但这次\textbf{染料有精确配方}。

\subsection{配方：$\varepsilon$-对齐会消去 $\varepsilon$，也会"晒出"卷积的形状}

$h[j'] = (y_2[j'] - \varepsilon a_0[j'])/\gamma_2$，其中 $a_0[j'] = \sum_t(\pm)R[(j'-p_t)]$
不带 $\varepsilon$。乘 $\hat\varepsilon$ 后 $a_0$ 贡献不带 $\varepsilon$ 地进入期望，
而
\[
\mathbb{E}[a_0[j]] = \frac{\tau}{d}\bigl(2\,\mathrm{CumR}[j] - \Sigma R\bigr)
\]
——$R$ 的\textbf{前缀和}（随机游走）。实测逐位置场与该式相关 0.989。

\subsection{传播到 $\Phi$：tent 核}

\[
\mathbb{E}[\Phi_R[i][m]] = -\frac{1}{\gamma_2}\Bigl[
s M R_i[m] + \beta\sum_v K[m,v]R_i[v]
+ s P_0 M (1-\tfrac{2m}{d})\mathbf{1}[i=0]\Bigr],
\]
$K[m,v] = d - 2|m-v|$（tent 核），$\beta = N\tau(\tau-1)/d^2$。
位置互斥 Monte-Carlo 算子与解析 tent 相关 0.9911。

\begin{intuition}[本章的"啊哈"]
$z_1$ 的确定性结构需要建模才能消除；hint 的结构\textbf{写得出解析式}——
系数只含公开常数（$M,\beta,K,d,\gamma_2,P_0$）。
\textbf{写得出解析式，就能求逆}。
\end{intuition}

\section{从 $R$ 到等价私钥}

\subsection{算子求逆 + 逐行标度}

解 $(sM\mathbb{I} + N F_{\mathrm{op}})\hat R_i = -\gamma_2\Phi_R[i] + $ 行 0 修正。
$s$ 逐行用整数聚类目标 $\frac1d\sum_m\cos(2\pi\hat R_i[m])$ 拟合
（真值处聚于整数，目标 $\approx0.8$）。

\subsection{mod-$q$ 线性求解}

$A_0 s_0 \equiv 16b_1 + R \pmod q$：768 个方程、768 个未知数，$A_0$ 模素数可逆，
解唯一精确。

\subsection{nbad razor：一把锋利的判刀}

$s_0$ 错一个坐标会通过 $A_0^{-1}$ \textbf{稠密扩散}到全部 768 坐标。
$\mathrm{nbad} = \#\{j: s_0[j]\notin\{-1,0,1\}\}$：真修正 $\Rightarrow 0$；
错误修正 $\Rightarrow \approx768$。收尾器只接受 $\mathrm{nbad}=0$ 的修正
（教训：贪心接受"nbad 减一"会毁掉全局）。

\subsection{等价私钥}

$\hat e_m = -\mathbf{1}[R_m\ge8] + \mathbf{1}[R_m\le-9]$，$\hat b_0 = R+\hat e$，
则 $A_0 s_0 + \hat e \equiv 16b_1 + \hat b_0$——$(s_0,\hat e,\hat b_0)$ 与真钥
签名能力相同。实测 $\hat e$ 非零 26 坐标（理论预期 $\approx32$，吻合）。

\section{实验时间线：每一步回答了什么问题}

\begin{story}[第一天：基础设施与"染色"发现]
\textbf{阶段 A--C}：公开数据管线与 C 内部值逐字节一致（10$^6$/10$^6$）；
$z_0$ 标签 96.19\%；$z_1$ 统计量三臂对照建立。
\textbf{阶段 5}：同键双数据集 $\Phi$ 零位置模式相关 0.9907——确定性结构存在。
\textbf{阶段 6/6B}：ICM 精化发散；逐坐标去场失败；场的来源定位为选择耦合。
\textbf{阶段 E/F}：代数上等价私钥成立（$R=b_0-e$ 逐系数验证）；oracle 密钥
20/20 通过原始验证器——机制通，当时误以为攻击路径未通（见下）。
\end{story}

\begin{story}[第二天：文献、hint 与闭合]
\textbf{上午}：文献定位与代码审计——重新读 sign.c 找未被利用的公开量。
\textbf{hint 审计}：发现 $y_2$ 抵消 $\Rightarrow$ hint 是 $R$ 的直接观测。
\textbf{tier-1 验证}：MC 斜率、池化信号、逐位置场——通道真实。
\textbf{累加与意外}：$10^6$ 累加直接分类失败——决定\textbf{推导}那个结构。
\textbf{场定律}：tent 算子与实测 corr 0.997。
\textbf{求逆与收尾}：oracle 臂单修正闭合 $\to$ $s_0$ 768/768；attack 臂 7 错
（iid 噪声跨界）。
\textbf{扩样闭合}：合并 $4\times10^6$ 签名，逐行标度，\textbf{直接零错误}，
$s_0$ 768/768；\textbf{伪造 20/20}。
\textbf{交叉验证}：七项检查全部通过。
\end{story}

\begin{story}[第三天（2026-09-28）：更正与复现]
我们在自己的数据与实现上\textbf{完整验证}了 $z_1$ 加权回归路线：
RMSE 0.184@20万、直接错误 6/768、公钥修正后闭合——恢复出\textbf{同一把}等价私钥。
这同时带来一个重要更正：我们此前"$z_1$ 路线饱和、不收敛"的判断是\textbf{错误}的
（那是我们早期压缩统计量与阈值判决的缺陷，不是通道的缺陷）。$z_1$ 回归路线成立；
hint 通道是第二条独立路线——两路线互为独立见证。
\end{story}

\section{交叉验证：凭什么信}

\begin{enumerate}
\item 不同代码路径：FFT $[x,-x]$ 独立重算负循环积，$R=b_0-e$ 精确。
\item 跨轮次：三份密钥文件（阶段 E / hint-oracle / hint-attack）SHA-256 一致。
\item 不同数据组合：real+sim1m（不含第三套）独立重跑——同样零错误、同一密钥。
\item 跨数据集：两套数据独立拟合的 $\hat R$ 一致 761/768，corr 0.999。
\item 实现无关性：累加器两种分块逐元素一致；快速版与慢速版逐元素一致。
\item 原始 C 全量复核：$3\times10^6$ 条签名全部通过未修改的原始验证器；
伪造 20/20 通过、4 组阴性对照全部拒绝。
\item 第二条路线交叉验证：$z_1$ 回归路线在同一数据上闭合，恢复同一把等价私钥。
\end{enumerate}

\begin{intuition}
任何单一实验都可能出错；但"换代码、换数据、换轮次、换验证器、甚至换一条
攻击路线"之后\textbf{同一个结论反复出现}，出错的可能性就被压到很低。
\end{intuition}

\section{常见问题（室友友好版 FAQ）}

\textbf{Q1：为什么不直接暴力枚举私钥？}
$s_0$ 有 $3^{768}\approx10^{366}$ 种取值。可观测宇宙的原子数约 $10^{80}$。
暴力在数学上不可行——所以才有统计攻击这种迂回打法。

\textbf{Q2：等价私钥是真私钥吗？}
不是。攻击者拿到的是 $(s_0, \hat e, \hat b_0)$——满足同一条代数关系、签名能力相同、
但个别坐标与真钥不同。对验证器二者不可区分——危害同级。

\textbf{Q3：会不会只是统计巧合？}
恢复出的 $s_0$ 与真钥 768/768 逐系数一致；两条独立路线恢复出同一把钥；
在未修改的原始 C 签名器上 20/20 通过原始验证器，4 组对照全部拒绝。巧合无法伪造这些。

\textbf{Q4：需要量子计算机吗？}
不需要。纯经典统计攻击：读签名、算直方图、解线性方程。笔记本就能跑。

\textbf{Q5：修复难吗？}
方向明确：不共享符号位、加宽 $z_0$ 掩码、重审 hint 论证式。工程量不大，
但安全证明要按混合分布模型重做。

\textbf{Q6：为什么叫"双峰"？}
响应的目标分布有两个对称的峰（来自 $\pm\varepsilon$），像骆驼双峰——
这也是掩码可减半的原因。

\textbf{Q7：$4\times10^6$ 签名现实吗？}
签名是公开量：一个服务每秒签几千条，一两天就攒够。攻击管线单机数小时跑完。
（$z_1$ 回归路线按其报告声称 20--25 万条即可；两条路线互补。）

\textbf{Q8：你们的工作和论坛上那份工作的关系？}
同一病根（共享符号位 + 联合分布泄露）、同一目标（等价私钥 + 伪造）。
恢复路线不同：他们用 $z_1$ 回归 + 置信窗口 Gray 码修正；我们用 hint 通道 +
场算子求逆 + 全解一致性 razor。我们复现了他们的路线（独立见证），
他们没有我们的 hint 通道（其论文只列为未来方向）。

\section{符号速查表}

\begin{center}
\begin{tabular}{@{}ll@{}}
\toprule
符号 & 含义 \\
\midrule
$d=256,\ q=26881$ & 维度与模数；$x^d\equiv-1$ \\
$\tau=30$ & 挑战非零系数个数（全为 $+1$） \\
$s_0, e$ & 主私钥与误差（短向量） \\
$b_1$（公开）/ $b_0$（保密） & $b=A_0s_0+e=16b_1+b_0$ 的高/低位 \\
$R=b_0-e$ & 攻击目标（公钥可计算） \\
$y_0,y_1,y_2$ & 掩码（3/10/11-bit 三角） \\
$z_0,z_1,z_2$ & 响应 $=$ 掩码 $+$ 秘密卷积 \\
$\varepsilon=(-1)^b$ & 共享双峰符号位 \\
$\Lambda$ & $z_0$ 通道的对数似然比（96.2\% 标签） \\
$\mathrm{cell},h$ & $\gamma_2$-格号与 hint \\
$\Phi_R[i][m]$ & hint 通道的对齐累加统计量 \\
$K[m,v]=d-2|m-v|$ & 场的 tent 核 \\
$\mathrm{nbad}$ & 一致性判刀 \\
\bottomrule
\end{tabular}
\end{center}

\vspace{0.6em}
\begin{center}
{\small 全部数字可复现：脚本、种子、产物见技术报告；\\
配套材料：\texttt{BiT128\_attack\_report.pdf}（严谨版）、\texttt{BiT128\_seminar\_beamers.pdf}（讨论班）。}
\end{center}

\end{document}
